\paragraph{Gödel--Zeckendorf 地址临界测度、负谱单支障碍与原子激活正规形}
在 Gödel--Zeckendorf 地址极限的严格自仿射闭合、纯碰撞位势的 kernel RH 临界律、primitive 原子剥离以及激活分支的局部解析展开已经建立之后，还可进一步抽取出下列七条彼此衔接的结构后果。第一条把黄金地址极限的临界几何常数完全闭式化；中间四条把 RH 阈值、解析缓冲层、两步原子扇区及其 Abel 有限部校正分别压缩为单支障碍、显式长度、一维 Witt 校正与闭式常数位移；最后两条则把激活分支在 \(u=1\) 附近的局部坐标先提升为逆正规形，再进一步压缩为无二次项的二阶沉默。

\begin{theorem}[Gödel--Zeckendorf 地址极限的临界 Hausdorff 测度闭式]\label{thm:xi-time-part62d-zg-critical-hausdorff-measure}
\leanverified{paper\_xi\_time\_part62d\_zg\_critical\_hausdorff\_measure}
沿用定理 \ref{thm:xi-zg-address-binary-self-affine}、\ref{thm:xi-time-part9g-holographic-prefix-isometry-on-line} 与
\ref{thm:xi-zg-address-minkowski-dimension-content} 的记号。记
\[
s_B:=\frac{\log\varphi}{\log B}.
\]
则 \(K_{B,v}\) 的 \(s_B\)-维 Hausdorff 测度满足
\[
\mathcal H^{s_B}(K_{B,v})=\frac{\varphi^2}{\sqrt5}.
\]
\end{theorem}
\begin{proof}
由定理 \ref{thm:xi-zg-address-minkowski-dimension-content}，对每个 \(n\ge 0\)，覆盖 \(K_{B,v}\) 所需的最少 \(B^{-n}\)-球个数恰为
\[
N_{B,v}(n)=F_{n+2}.
\]
因此用这组球直接覆盖可得
\[
\mathcal H_{B^{-n}}^{s_B}(K_{B,v})
\le
F_{n+2}B^{-ns_B}.
\]

另一方面，由定理 \ref{thm:xi-time-part9g-holographic-prefix-isometry-on-line} 的柱--球对应，长度 \(n\) 的不同 admissible 前缀落在彼此不同的 \(B^nL_v\)-余类中。故任一直径不超过 \(B^{-n}\) 的集合都至多与其中一个余类相交，从而任意这类覆盖都至少需要 \(F_{n+2}\) 个集合。于是
\[
\mathcal H_{B^{-n}}^{s_B}(K_{B,v})
\ge
F_{n+2}B^{-ns_B}.
\]
故对每个 \(n\) 都有
\[
\mathcal H_{B^{-n}}^{s_B}(K_{B,v})
=
F_{n+2}B^{-ns_B}.
\]

再由 \(B^{s_B}=\varphi\) 与 Binet 公式，
\[
F_{n+2}B^{-ns_B}=F_{n+2}\varphi^{-n}\longrightarrow \frac{\varphi^2}{\sqrt5}.
\]
当 \(n\to\infty\) 时 \(B^{-n}\downarrow 0\)，故由 Hausdorff 外测度的单调极限得到
\[
\mathcal H^{s_B}(K_{B,v})
=
\lim_{n\to\infty}\mathcal H_{B^{-n}}^{s_B}(K_{B,v})
=
\frac{\varphi^2}{\sqrt5}.
\]
证毕。
\end{proof}

\begin{theorem}[核版 RH/Ramanujan 障碍由唯一负谱支完全控制]\label{thm:xi-time-part62d-collision-rh-negative-branch-control}
沿用命题 \ref{prop:real-input-40-collision-all-real} 与
\ref{prop:real-input-40-collision-rhk-window} 的记号。设
\[
f(\lambda;u)=\lambda^3-2\lambda^2-(u+1)\lambda+u
\]
的三条实根按
\[
\lambda_+(u)>1,\qquad 0<\lambda_0(u)<1,\qquad \lambda_-(u)<0
\]
排序，并记
\[
\Lambda(u):=\max\{|\lambda_0(u)|,|\lambda_-(u)|,1\}.
\]
则 kernel RH 条件等价于
\[
\textsf{RH}(u)\iff
\begin{cases}
\text{恒真}, & 0<u\le 1,\\[1mm]
(-\lambda_-(u))^2\le \lambda_+(u), & u>1.
\end{cases}
\]
特别地，\(\lambda_0(u)\) 从不决定 RH 窗口边界；唯一的障碍分支始终是负根 \(\lambda_-(u)\)。
\end{theorem}
\begin{proof}
由命题 \ref{prop:real-input-40-collision-all-real}，
\[
0<\lambda_0(u)<1
\qquad(\forall\,u>0),
\]
故 \(|\lambda_0(u)|<1\) 对全部正参数恒成立。

当 \(0<u\le 1\) 时，同一命题给出
\[
\lambda_-(u)\in[-1,0),
\]
于是
\[
|\lambda_-(u)|\le 1,\qquad |\lambda_0(u)|<1.
\]
从而
\[
\Lambda(u)=1.
\]
又因 \(\lambda_+(u)>1\)，故
\[
1\le \lambda_+(u)^{1/2},
\]
于是 \(\textsf{RH}(u)\) 自动成立。

当 \(u>1\) 时，命题 \ref{prop:real-input-40-collision-all-real} 给出
\[
\lambda_-(u)<-1.
\]
因此
\[
|\lambda_-(u)|>1>|\lambda_0(u)|,
\]
从而
\[
\Lambda(u)=|\lambda_-(u)|=-\lambda_-(u).
\]
代回命题 \ref{prop:real-input-40-collision-rhk-window} 的判据
\[
\textsf{RH}(u)\iff \Lambda(u)\le \lambda_+(u)^{1/2}
\]
便得
\[
\textsf{RH}(u)\iff -\lambda_-(u)\le \lambda_+(u)^{1/2}
\iff
(-\lambda_-(u))^2\le \lambda_+(u).
\]
证毕。
\end{proof}

\begin{proposition}[RH 临界参数的结果式消元范式]\label{prop:xi-time-part62d-collision-rh-resultant-template}
设
\[
f(\lambda;u)\in \RR[\lambda,u]
\]
是关于 \(\lambda\) 的首一三次多项式，并且在某个区间 \(I\subset(0,\infty)\) 上具有三条实根分支
\[
\lambda_+(u)>\lambda_0(u)>\lambda_-(u),
\qquad
\lambda_+(u)>0>\lambda_-(u).
\]
若在某个 \(u\in I\) 处发生 Ramanujan 型达界事件
\[
|\lambda_-(u)|=\sqrt{\lambda_+(u)},
\]
则必有
\[
\operatorname{Res}_{\lambda}\bigl(f(\lambda;u),f(\lambda^2;u)\bigr)=0.
\]
对纯碰撞三次
\[
f(\lambda;u)=\lambda^3-2\lambda^2-(u+1)\lambda+u
\]
应用该范式，可得精确因式分解
\[
\operatorname{Res}_{\lambda}\bigl(f(\lambda;u),f(\lambda^2;u)\bigr)
=
2u\bigl(u^5+4u^4+3u^3-96u^2-42u-14\bigr).
\]
因此正参数域中的 Ramanujan 临界候选恰由五次方程
\[
u^5+4u^4+3u^3-96u^2-42u-14=0
\]
给出。
\end{proposition}
\begin{proof}
若 \(|\lambda_-(u)|=\sqrt{\lambda_+(u)}\)，由于 \(\lambda_-(u)<0\)，便有
\[
\lambda_-(u)^2=\lambda_+(u).
\]
取 \(\lambda=\lambda_-(u)\)，则
\[
f(\lambda;u)=0,
\qquad
f(\lambda^2;u)=f(\lambda_+(u);u)=0.
\]
故 \(f(\lambda;u)\) 与 \(f(\lambda^2;u)\) 在变量 \(\lambda\) 上有公共根，于是其结果式必为零。

对给定三次
\[
f(\lambda;u)=\lambda^3-2\lambda^2-(u+1)\lambda+u,
\]
直接计算 Sylvester 结果式即得
\[
\operatorname{Res}_{\lambda}\bigl(f(\lambda;u),f(\lambda^2;u)\bigr)
=
2u\bigl(u^5+4u^4+3u^3-96u^2-42u-14\bigr).
\]
由此可知所有正参数临界候选都必须落在该五次方程的正根集合中。再结合命题 \ref{prop:real-input-40-collision-rhk-window} 的唯一性结论，便得到真正的临界参数 \(u_R\)。证毕。
\end{proof}

\begin{corollary}[解析但非 RH 实轴相的显式长度与负谱触发机制]\label{cor:xi-time-part62d-analytic-nonrh-window-length}
沿用定理 \ref{thm:xi-rh-critical-before-nearest-complex-branch} 与
\ref{thm:xi-time-part62d-collision-rh-negative-branch-control} 的记号。记
\[
\delta_{\mathrm{an}}:=R_\theta-\theta_R.
\]
则
\[
\delta_{\mathrm{an}}
=
R_\theta-\theta_R
\approx
1.48342289346087>0.
\]
并且对任意
\[
\theta\in(\theta_R,R_\theta)
\]
都有
\[
(-\lambda_-(e^\theta))^2>\lambda_+(e^\theta).
\]
换言之，在这一非空实区间内，Perron 压力分支仍保持解析，而 kernel RH 的失效已经由负谱支的平方越界单独触发。
\end{corollary}
\begin{proof}
由定理 \ref{thm:xi-rh-critical-before-nearest-complex-branch}，
\[
\theta_R<R_\theta,
\qquad
R_\theta\approx 2.76230289346087,
\qquad
\theta_R\approx 1.27888.
\]
故
\[
\delta_{\mathrm{an}}=R_\theta-\theta_R\approx 1.48342289346087.
\]

若 \(\theta\in(\theta_R,R_\theta)\)，则首先 \(|\theta|<R_\theta\)，因而 Perron 分支 \(P_2(\theta)\) 仍位于解析收敛盘内部。另一方面 \(u=e^\theta>\!u_R>1\)，由定理 \ref{thm:xi-time-part62d-collision-rh-negative-branch-control}，
\[
\textsf{RH}(u)\iff (-\lambda_-(u))^2\le \lambda_+(u).
\]
而此时 RH 已失效，故必有
\[
(-\lambda_-(u))^2>\lambda_+(u).
\]
把 \(u=e^\theta\) 代回即得所示不等式。证毕。
\end{proof}

\begin{theorem}[两步原子扇区在 primitive 层的一维 Witt 塌缩]\label{thm:xi-time-part62d-atomic-two-step-witt-rankone-collapse}
沿用命题 \ref{prop:atomic-2-orbit-split-logM} 与定理
\ref{thm:xi-atomic-witt-cycle-primitive-exact-peeling} 的记号。记
\[
P(z,u)=\sum_{n\ge 1}p_n(u)z^n,
\qquad
P_{\mathrm{core}}(z,u)=\sum_{n\ge 1}p_n^{\mathrm{core}}(u)z^n.
\]
则：
\begin{enumerate}
  \item 对一切 \(n\neq 2\)，有
  \[
  p_n(u)=p_n^{\mathrm{core}}(u);
  \]
  \item 在 \(n=2\) 处唯一出现原子修正
  \[
  p_2(u)=u+p_2^{\mathrm{core}}(u).
  \]
\end{enumerate}
此外，命题 \ref{prop:atomic-2-orbit-split-logM} 所嵌审计输出表明，branch 层满足
\[
(|\tau|,E)=(2,1)
\]
的两步回路恰有 \(7\) 条。因而这 \(7\) 条几何上彼此不同的两步 branch 回路在 primitive--Witt 层的净效应严格压缩为单一 monomial
\[
u z^2.
\]
等价地，
\[
\dim_{\QQ(u)}
\operatorname{span}\{\,p_n(u)-p_n^{\mathrm{core}}(u):n\ge 1\,\}=1.
\]
\end{theorem}
\begin{proof}
定理 \ref{thm:xi-atomic-witt-cycle-primitive-exact-peeling} 已给出精确剥离
\[
P(z,u)=u z^2+P_{\mathrm{core}}(z,u).
\]
逐项比较系数，立即得到
\[
p_n(u)=p_n^{\mathrm{core}}(u)\quad(n\neq 2),
\qquad
p_2(u)=u+p_2^{\mathrm{core}}(u).
\]
故 full/core 的 primitive 差分仅支撑在长度 \(2\) 的单一坐标上，其 \(\QQ(u)\)-线性张成空间因此恰为一维。

另一方面，命题 \ref{prop:atomic-2-orbit-split-logM} 同页嵌入的审计输出已经逐项列出 branch 层全部满足
\[
(|\tau|,E)=(2,1)
\]
的两步回路，并给出其总数恰为 \(7\)。于是这 \(7\) 条 branch 两步回路在 primitive--Witt 层的总校正只能通过唯一项
\[
u z^2
\]
出现。证毕。
\end{proof}

\begin{theorem}[七条两步 branch 回路在 primitive--Witt 与 Abel 有限部中只留下单原子校正]\label{thm:xi-time-part62d-seven-two-step-branch-loops-single-atom-finitepart-shift}
沿用命题 \ref{prop:atomic-2-orbit-split-logM} 与定理
\ref{thm:xi-time-part62d-atomic-two-step-witt-rankone-collapse} 的记号。令
\[
P(z)=\sum_{n\ge 1}p_nz^n,
\qquad
P_{\mathrm{core}}(z)=\sum_{n\ge 1}p_n^{\mathrm{core}}z^n
\]
表示 \(u=1\) 时的 primitive 系列，并记
\[
P(z)=-\log(1-\lambda z)+\log\mathfrak M+o(1),
\]
\[
P_{\mathrm{core}}(z)=-\log(1-\lambda z)+\log\mathfrak M_{\mathrm{core}}+o(1)
\qquad
(z\to z_\star),
\]
其中
\[
\lambda=\varphi^2,
\qquad
z_\star=\lambda^{-1}=\varphi^{-2}.
\]
则
\[
p_2=1+p_2^{\mathrm{core}},
\qquad
p_n=p_n^{\mathrm{core}}
\quad(n\neq 2).
\]
特别地，branch 层中恰为 \(7\) 条的两步回路在 primitive--Witt 层只留下唯一原子
\[
z^2.
\]
并且 Abel 有限部常数满足精确位移
\[
\log\mathfrak M-\log\mathfrak M_{\mathrm{core}}
=
z_\star^2
=
\varphi^{-4}
=
\frac{7-3\sqrt5}{2}.
\]
等价地，
\[
\log\mathfrak M_{\mathrm{core}}
=
\log\mathfrak M-\varphi^{-4}.
\]
\end{theorem}
\begin{proof}
将定理 \ref{thm:xi-time-part62d-atomic-two-step-witt-rankone-collapse} 在 \(u=1\) 处专门化，即得
\[
P(z)=z^2+P_{\mathrm{core}}(z).
\]
逐项比较系数便得到
\[
p_2=1+p_2^{\mathrm{core}},
\qquad
p_n=p_n^{\mathrm{core}}
\quad(n\neq 2).
\]
同一定理已经说明，branch 层满足
\[
(|\tau|,E)=(2,1)
\]
的两步回路恰有 \(7\) 条，故它们在 primitive--Witt 层的净效应正是唯一原子 \(z^2\)。

另一方面，命题 \ref{prop:atomic-2-orbit-split-logM} 给出精确常数分解
\[
\log\mathfrak M
=
\log\mathfrak M_{\mathrm{core}}+z_\star^2.
\]
再代入
\[
z_\star=\varphi^{-2}
\]
即得
\[
\log\mathfrak M-\log\mathfrak M_{\mathrm{core}}
=
\varphi^{-4}
=
\frac{7-3\sqrt5}{2}.
\]
证毕。
\end{proof}

\begin{theorem}[激活分支在 \texorpdfstring{$u=1$}{u=1} 的逆正规形与无二次重整化]\label{thm:xi-time-part62d-activated-branch-inverse-normalform}
沿用注 \ref{rem:real-input-40-activated-branch-series} 与定理
\ref{thm:xi-output-potential-activated-branch-rigidity} 的记号。令
\[
h:=u-1,
\qquad
\lambda:=\lambda_{\mathrm{act}}(u).
\]
则正向展开
\[
\lambda
=
h+3h^3-h^4+18h^5-22h^6+O(h^7)
\]
在局部可逆，并且其反函数满足显式逆正规形
\[
h
=
\lambda-3\lambda^3+\lambda^4+9\lambda^5+\lambda^6+O(\lambda^7).
\]
特别地：
\begin{enumerate}
  \item 正向展开与逆正规形中都不存在二次项；
  \item 第一非线性修正项严格为三次项；
  \item 在逆正规形一侧，该三次系数恰为 \(-3\)。
\end{enumerate}
\end{theorem}
\begin{proof}
由注 \ref{rem:real-input-40-activated-branch-series}，
\[
\lambda
=
h+3h^3-h^4+18h^5-22h^6+O(h^7).
\]
设其局部反函数写成
\[
h
=
\lambda+a_2\lambda^2+a_3\lambda^3+a_4\lambda^4+a_5\lambda^5+a_6\lambda^6+O(\lambda^7).
\]
把此式代回上面的正向展开，并按 \(\lambda\) 的幂次逐项比较。

\(\lambda^2\) 项给出
\[
a_2=0.
\]
\(\lambda^3\) 项给出
\[
a_3+3=0,
\]
故
\[
a_3=-3.
\]
\(\lambda^4\) 项给出
\[
a_4-1=0,
\]
故
\[
a_4=1.
\]
\(\lambda^5\) 项给出
\[
a_5+9a_3+18=0,
\]
从而
\[
a_5=9.
\]
\(\lambda^6\) 项给出
\[
a_6+9a_4-4a_3-22=0,
\]
故
\[
a_6=1.
\]
于是
\[
h
=
\lambda-3\lambda^3+\lambda^4+9\lambda^5+\lambda^6+O(\lambda^7).
\]
所述三条推论随即成立。证毕。
\end{proof}

\begin{theorem}[激活支的二阶沉默]\label{thm:xi-time-part62d-activated-branch-second-order-silence}
沿用定理 \ref{thm:xi-time-part62d-activated-branch-inverse-normalform} 的记号，记
\[
h:=u-1.
\]
则激活分支满足显式展开
\[
\lambda_{\mathrm{act}}(u)
=
h+3h^3-h^4+18h^5-22h^6+O(h^7),
\]
从而
\[
\lambda_{\mathrm{act}}(u)-(u-1)=O\!\bigl((u-1)^3\bigr).
\]
特别地，激活支在 \(u=1\) 邻域内不存在二次弯曲项，首个非线性修正严格出现在三次项。
\end{theorem}
\begin{proof}
定理 \ref{thm:xi-time-part62d-activated-branch-inverse-normalform} 的正向展开已给出
\[
\lambda_{\mathrm{act}}(u)
=
h+3h^3-h^4+18h^5-22h^6+O(h^7).
\]
从该式直接读出二次项系数为 \(0\)，并把 \(h=u-1\) 代回，即得
\[
\lambda_{\mathrm{act}}(u)-(u-1)=O\!\bigl((u-1)^3\bigr).
\]
证毕。
\end{proof}

\noindent
上述七条结果把 H.160--H.162 的 \(B\)-进地址链与 D.509--D.532 的纯谱/Witt/激活链压缩为同一层二值刚性：一方面，Gödel--Zeckendorf 地址极限在临界维数处给出完全显式的黄金常数；另一方面，kernel RH 的破裂不再依赖多支比较，而被负谱单支完全控制，并在解析域内部留下显式长度的缓冲层；与此同时，两步原子扇区在 primitive 层只留下一个一维方向，而激活分支的局部坐标直到二阶都保持严格线性，首个非线性修正只能在三阶显影。于是，地址几何、谱阈值、Witt 剥离与局部正规形便被统一压缩为同一条可审计的刚性主线。

\endinput
